% Created 2025-12-02 mar 23:29
% Intended LaTeX compiler: pdflatex
\documentclass[11pt]{article}
\usepackage[utf8]{inputenc}
\usepackage[T1]{fontenc}
\usepackage{graphicx}
\usepackage{longtable}
\usepackage{wrapfig}
\usepackage{rotating}
\usepackage[normalem]{ulem}
\usepackage{amsmath}
\usepackage{amssymb}
\usepackage{capt-of}
\usepackage{hyperref}
\date{\today}
\title{Migracion}
\hypersetup{
 pdfauthor={},
 pdftitle={Migracion},
 pdfkeywords={},
 pdfsubject={},
 pdfcreator={Emacs 30.2 (Org mode 9.7.34)}, 
 pdflang={English}}
\begin{document}

\maketitle
\tableofcontents

\section{Migración Wordpress}
\label{sec:org134793f}

\subsection{Primeros Pasos}
\label{sec:org5bbb830}

\begin{enumerate}
\item Revisamos usuarios mysql

\begin{verbatim}
sudo mariadbysql -u root
\end{verbatim}

\begin{verbatim}
SELECT User, Host, Password, is_role FROM mysql.user;
\end{verbatim}
\end{enumerate}


\begin{enumerate}
\item Cambiamos passwords

\begin{verbatim}
ALTER USER 'root'@'localhost' IDENTIFIED BY 'passwd';
FLUSH PRIVILEGES;
\end{verbatim}
\end{enumerate}


\begin{enumerate}
\item Vemos los usuarios que hay

\begin{verbatim}
SELECT User, Host, Password, is_role FROM mysql.user;
\end{verbatim}
\end{enumerate}
\subsection{Wordpress a Debian 13}
\label{sec:org989110c}

\begin{enumerate}
\item Primero necesitamos el sql de la base de datos a migrar

\begin{verbatim}

mysqldump -h HOST_DB -u USUARIO_DB -p NOMBRE_DB > dump.sql

# opcional --no-tablespaces

\end{verbatim}

\item Ahora es necesario modificar el sql

\begin{itemize}
\item Será bueno hacer un backup del sql original
\end{itemize}

\begin{verbatim}
VIEJO="dominio_antiguo"
NUEVO="dominio_nuevo"

sed -i "s/$VIEJO/$NUEVO/g" dump.sql
sed -i '/^USE /d' dump.sql
sed -i '/^CREATE DATABASE /d' dump.sql
\end{verbatim}

\item Creamos directorios para la migración

\begin{verbatim}
mkdir -p /var/www/html/directorio
\end{verbatim}

\item Copiamos los archivos y les damos los permisos necesarios

\begin{verbatim}

# Copiamos el dump de la base de datos
scp user@host:/ruta/dump.sql /ruta/nueva/

# Copiamos wordpress entero
scp user@host:/ruta /var/www/html/directorio

sudo chown -R www-data:www-data /var/www/html/directorio
chmod -R 755 /var/www/html/directorio
\end{verbatim}

\item Creamos la base de datos y usuarios con permisos necesarios

\begin{verbatim}
sudo mariadb -u root -p
\end{verbatim}

\begin{verbatim}
CREATE DATABASE nueva_db;
CREATE USER 'nuevo_user'@'localhost' IDENTIFIED BY 'passwd';

GRANT ALL PRIVILEGES ON nueva_db.* TO 'nuevo_user'@'localhost';
FLUSH PRIVILEGES;
\end{verbatim}

\begin{itemize}
\item Importamos la base de datos a la nueva es importante que el sql esté modificado con los dominios nuevos

\begin{verbatim}
mysql -u nuevo_user -p nueva_db < dump.sql
\end{verbatim}

\item Comprobamos que existan las tablas, dentro de mariadb

\begin{verbatim}
USE nueva_db;
SHOW TABLES;
\end{verbatim}
\end{itemize}

\item Ahora hay que conectar la base de datos al wordpress

\begin{itemize}
\item Hay que actualizar el archivo wp-config.php
\end{itemize}

\begin{verbatim}
define( 'DB_NAME', 'nueva_db' );
define( 'DB_USER', 'nuevo_user' );
define( 'DB_PASSWORD', 'passwd' );
define( 'DB_HOST', 'localhost' );
\end{verbatim}

\item Modificamos el archivo .htaccess

\begin{itemize}
\item Podemos entrar al administrador de wordpress, en wp-admin.
\begin{itemize}
\item Entrar en Ajustes → Enlaces Permanentes
\item Guardar cambios (2 veces)
\end{itemize}

\item Manualmente:

\begin{verbatim}
<IfModule mod_rewrite.c>
    RewriteEngine On
    RewriteBase /picante/
    RewriteRule ^index\.php$ - [L]
    RewriteCond %{REQUEST_FILENAME} !-f
    RewriteCond %{REQUEST_FILENAME} !-d
    RewriteRule . /picante/index.php [L]
</IfModule>
# END WordPress
\end{verbatim}

\begin{quote}
\emph{picante} es la ruta donde está instalado el wordpress sería \emph{var/www/html/picante}
\end{quote}
\end{itemize}

\item Activar mod\textsubscript{rewrite}

\begin{itemize}
\item Habilitamos el módulo de reescritura de URLs de Apache

\begin{verbatim}
sudo a2enmod rewrite

sudo systemctl restart apache2
\end{verbatim}
\end{itemize}

\item Allow override All en el directory de Apache

\begin{itemize}
\item Hay que permitir que Apache lea tu fichero .htaccess

\item Editamos /etc/apache2/sites-available/000-default.conf

\begin{itemize}
\item Añadimos el bloque <Directory> dentro de <VirtualHost *:80>
\end{itemize}

\begin{verbatim}
# ... otras configuraciones ...
DocumentRoot /var/www/html
# AÑADE ESTE BLOQUE:
    <Directory /var/www/html/mi-web>
        Options Indexes FollowSymLinks
        AllowOverride All
        Require all granted
    </Directory>
</VirtualHost>
\end{verbatim}

\item Ejemplo:
\begin{verbatim}
<VirtualHost *:80>
	# The ServerName directive sets the request scheme, hostname and port that
	# the server uses to identify itself. This is used when creating
	# redirection URLs. In the context of virtual hosts, the ServerName
	# specifies what hostname must appear in the request's Host: header to
	# match this virtual host. For the default virtual host (this file) this
	# value is not decisive as it is used as a last resort host regardless.
	# However, you must set it for any further virtual host explicitly.
	#ServerName www.example.com

	ServerAdmin webmaster@localhost
	DocumentRoot /var/www/html

	# Available loglevels: trace8, ..., trace1, debug, info, notice, warn,
	# error, crit, alert, emerg.
	# It is also possible to configure the loglevel for particular
	# modules, e.g.
	#LogLevel info ssl:warn

	ErrorLog ${APACHE_LOG_DIR}/error.log
	CustomLog ${APACHE_LOG_DIR}/access.log combined

	# For most configuration files from conf-available/, which are
	# enabled or disabled at a global level, it is possible to
	# include a line for only one particular virtual host. For example the
	# following line enables the CGI configuration for this host only
	# after it has been globally disabled with "a2disconf".
	#Include conf-available/serve-cgi-bin.conf
	<Directory /var/www/html/picante>
		Options Indexes FollowSymLinks
		AllowOverride All
		Require all granted
	</Directory>

</VirtualHost>


\end{verbatim}
\end{itemize}

\item Reiniciar

\begin{verbatim}
sudo systemctl restart apache2
\end{verbatim}
\end{enumerate}
\end{document}
